\documentclass[10pt,a4paper]{amsart}
\usepackage[margin=19mm]{geometry}
\usepackage{amsmath,amssymb,mathtools}
\usepackage[scr=rsfs,cal=euler]{mathalfa}
\usepackage{xfrac}
\usepackage[unicode,hidelinks,bookmarks=false]{hyperref}
\newcommand{\ie}{i.e.\ }
\newcommand{\cf}{cf.\ }
\newcommand{\resp}{respectively\ }
\newcommand{\Subsection}{\subsection}
\providecommand{\nobreakdash}{\nobreak\hbox{-}}
\setlength{\emergencystretch}{2em}
\multlinegap0pt
\usepackage{bm}
\usepackage{bbm}
\usepackage{dsfont}
\usepackage{xspace}
\usepackage{cancel}
\usepackage{mathtools}
\usepackage{amsthm}
\makeatletter
\@ifundefined{theorem}{\newtheorem{theorem}{Theorem}}{}
\@ifundefined{lemma}{\newtheorem{lemma}[theorem]{Lemma}}{}
\makeatother
\newcommand\rounddown[1]{\left\lfloor#1\right\rfloor}
\newcommand\roundup[1]{\left\lceil#1\right\rceil}
\title[On graphs with girth at least five achieving Steffen's edge ]{On graphs with girth at least five achieving Steffen's edge coloring bound}
\author{Guantao Chen, Alireza Fiujlaali, Anna Johnsen-Yu, Jessica McDonald}
\date{Source: arXiv:2601.23274v1 (CC BY 4.0)}
\begin{document}
\maketitle

\noindent\emph{Source and licence.} On graphs with girth at least five achieving Steffen's edge coloring bound. Version arXiv:2601.23274v1 (CC BY 4.0), distributed under the Creative Commons Attribution 4.0 International licence, as recorded in the arXiv licence metadata for this exact version and shown on the abstract page https://arxiv.org/abs/2601.23274v1. Full rights notice: licenses/arxiv-2601.23274-CC-BY-4.0.txt.\par\smallskip

\noindent\textit{Editorial scope.} Counted unit: the Lemma whose source label is \texttt{critical} in the article (source0.tex lines 163--169, source label \texttt{critical}), delivered together with the article's own complete original proof of it (source0.tex lines 172--210, one \texttt{proof} environment of 39 lines). No mathematics is written, repaired or re-derived: statement and proof are the article's own text line for line. Required context retained uncounted: Thm-CJZ (lines 157--158). Every definition, display and label of those lines is retained unchanged. Omitted from this package and not redistributed: 1--156, 159--162, 170--171, 211--622. Nothing inside a retained excerpt is omitted or altered: every retained body is delivered exactly as the article's lines are written, so no omission receipt of any kind accompanies this package. Citations: the retained excerpts carry no citation command at all, so no bibliography entry of the article is transcribed and no reference is redistributed. Reference adaptation: none was needed and none was made; every reference inside the delivered unit resolves inside it, so no unresolved reference or citation is produced. Printed slips kept literal: no printed word, symbol, hypothesis or number of the counted statement or of its proof was altered. Layout and class: the article's own class is \texttt{amsart} with packages including inputenc, bm, verbatim, graphicx, subcaption, amssymb, color, mathabx, mathtools, enumitem, mathrsfs, amsmath, amsthm, ulem; the wrapper is the lane's pinned \texttt{amsart} editorial wrapper at 10pt on A4, which restates the theorem-like environments the retained text needs and adds the article's own macro definitions that the retained text needs (\texttt{\textbackslash rounddown}, \texttt{\textbackslash roundup}), with extra wrapper packages (bm, bbm, dsfont, xspace, cancel, mathtools). The delivered numbering is the unit's own. Assets: no retained span contains a figure environment, a graphics-inclusion command, a PDF-page inclusion, a table or a TikZ picture; no image file is copied into this package, the package declares no asset dependency and no image file is redistributed with it. Licence: version arXiv:2601.23274v1 of this article is distributed under the Creative Commons Attribution 4.0 International licence, as recorded in the arXiv licence metadata for this exact version and shown on the abstract page https://arxiv.org/abs/2601.23274v1. A full rights notice is delivered at licenses/arxiv-2601.23274-CC-BY-4.0.txt. Characters: every retained excerpt is pure ASCII, so the delivered engine, XeLaTeX with its default Latin Modern text font, reports no missing character.
\begin{theorem}[Chen, Hao, Jing, Yu, Zang]\label{Thm-CJZ}
For any graph $G$, if $\chi'(G) \ge \Delta(G) +2$ then $\chi'(G) = \Gamma(G)$. \end{theorem}

\begin{lemma} \label{critical}
    Let $G$ be a critical graph on $n$ vertices. If  $\chi'(G) \geq \Delta + 2$,
    then $n$ is odd and for every vertex $v\in V(G)$,
    $$d_{G}(v) = \sum_{w\ne v} (\chi' -1 -d_G(w)) +2  \ge (n-1)(\chi' -1 -\Delta) +2.$$
Moreover, if $n \ge g$ and $\chi'(G) = \Delta + \lceil \mu/\lfloor g/2 \rfloor\rceil \ge \Delta +2$ for some positive integer $g \ge 5$, then
    \( \delta(G) \ge  n\mu/g +1\), and  equality holds only if $\mu =g$.
    \end{lemma}

\begin{proof} By definition, $\Gamma(G) \le \chi'(G)$. Applying  Theorem~\ref{Thm-CJZ}, we find that $\Gamma(G)=\chi'(G)$ as $\chi'(G) \ge \Delta +2$. Since $G$ is critical, for every proper subgraph $H$ of $G$ with an odd number of vertices, we have that $\roundup{2|E(H)|/(|H|-1)} \le \Gamma(H) \le \chi'(H) < \chi'(G) = \Gamma(G)$.  Hence, $G$ itself is the only subgraph $H$ of $G$ with an odd number of vertices such that $\roundup{2|E(H)|/(|H|-1)} = \Gamma(G)$. Consequently, $n$ is odd, and for any edge $e\in E(G)$,
\( \chi'(G-e) =\chi'(G) -1\). Hence,
\[
 \roundup{\frac{|E(G-e)|}{(n-1)/2}} \le \chi'(G-e)= \chi'(G) -1 = \roundup{\frac{|E(G)|}{(n-1)/2}} -1.
\]
Since $\frac{|E(G-e)|}{\rounddown{(n-1)/2}} = \frac{|E(G)|-1}{(n-1)/2} \ge \roundup{\frac{|E(G)|}{(n-1)/2}} -1$, it follows that  $\frac{|E(G-e)|}{\rounddown{(n-1)/2}}$ is an integer and is equal to $\chi'(G-e)$. Hence, for an edge-coloring of $G-e$ with $\chi'(G-e)$ colors, each color class contains exactly $(n-1)/2$ edges. Therefore, \(E(G - e)\) can be partitioned into \(\chi'(G) - 1\) near-perfect matchings, each of which misses exactly one vertex of \(G\).

Let $v$ be the  vertex specified in Lemma~\ref{critical}, and  let $e$ be an edge not incident with $v$.  For each vertex $w\ne v$, if $w$ is not incident to $e$, there are exactly $\chi'(G) -1 -d_G(w)$ of these near-perfect matchings not containing the vertex $w$;  if $w$ is incident to $e$, there are exactly $\chi'(G) -1 -(d_G(w) -1)$ of these matchings  not containing the vertex $w$. Note that every one of these matchings contains $v$. Hence, $d_G(v)  = \sum_{w\ne v} (\chi'(G) -1 -d_G(w)) +2$. As $d_G(w) \le \Delta(G)$, we have the following.
\[
\delta(G) \ge (n-1)(\chi'(G) -1 -\Delta(G)) +2.
\]

Now suppose $\chi'(G)=\Delta + \roundup{\mu/\rounddown{g/2}}\geq \Delta + 2$ for some positive integer $g\geq 5$. Then by substituting $\Delta + \roundup{\mu/\rounddown{g/2}}$ for $\chi'(G)$, we obtain the following inequality.
\[
\delta(G) \ge (n-1)\left(\roundup{\frac{\mu}{\rounddown{g/2}}}-1\right) +2\ge (n-1)\left(\roundup{{\frac{2\mu}g}}-1\right) +2.
\]

Note that
\[
(n-1)\left(\roundup{\frac{\mu}{\rounddown{g/2}}}-1\right) +2 -  \left(\frac {n\mu}g +1\right)
\ge (n-1)\left(\roundup{\frac{2\mu}g}-\frac {\mu}g -1\right) -\frac {\mu}g +1.
\]
If $\mu <   g$, then $\frac {\mu +1}g \le 1$, and hence
$\roundup{\frac{2\mu}g} \ge 2 \ge 1+ \frac {\mu +1}g$. Thus,
\[
(n-1)\left(\roundup{\frac{2\mu}g}-\frac {\mu}g -1\right) -\frac {\mu}g +1 \ge \frac 1g ((n-1) -\mu) +1 \ge \frac 1g((n-1) -g) +1> 0.
\]

If $\mu =g$, then
$
(n-1)\left(\roundup{\frac{2\mu}g}-\frac {\mu}g -1\right) -\frac {\mu}g +1= 0$, as desired.

We now turn to the case $\mu >g$. In this case, $\roundup{2\mu/g} \ge 3$. If $\mu/g \le 1.5$,  then $4\mu/(3g) \le 2$, and hence $\roundup{2\mu/g} =3\ge 1 +4\mu/(3g)$.  If $\mu/g > 1.5$, then $(2/3) (\mu/g) \ge 1$, which implies that $\mu/g > 1 + \mu/(3g)$, and so $\roundup{2\mu/g} \ge \mu/g +\mu/g \ge 1+\mu/3g + \mu/g = 1 +4\mu/(3g)$.
Hence,
\[
(n-1)\left(\roundup{\frac{2\mu}g}-\frac {\mu}g -1\right) -\frac {\mu}g +1\ge \frac{(n-1)\mu}{3g} -\frac{\mu}g +1 \ge \frac {4\mu} g -\frac {\mu}g +1  > 0,
\]
as $n \ge 5$ and $\mu/g \ge 1$.
\end{proof}

\end{document}
